% -------------------------------------------------------
\section{Abstract}
% -------------------------------------------------------

The evolution of scientific discourse is both an empirical question about a research field and a methodological test of text analytics. This study examines how AI, machine learning, and NLP research abstracts shifted from theoretical to applied orientations, changed their framing of limitations, and accelerated the spread of new terminology over the last three decades. Five complementary methods, lexical keyword analysis, topic modelling (LDA and NMF), limitation framing classification, terminology diffusion tracking (TF-IDF delta, Word2Vec, and SPECTER2), and dimensionality reduction (PCA and UMAP), are applied to 24,677 arXiv abstracts. All five methods independently identify 2014–2016 as the primary inflection point, with TF-IDF and Sentence-BERT (\texttt{all-MiniLM-L6-v2}) centroid distance series producing a Pearson correlation of $r = 0.863$. Application-oriented language overtook theory-oriented language around 2014–2015, ethics-related terms grew from near-zero to approximately 13\% of abstracts by 2020–2024, and UMAP suggests the field expanded into new paradigms rather than displacing existing ones. At the same time, within-year dispersion decreases significantly (Kendall’s $\tau \approx -0.5$), indicating increasing internal coherence as the field evolves. Most notably, the dominant mode of terminology change is quiet conceptual repositioning rather than frequency growth; 13 of 36 tracked terms shifted meaning without a proportionate rise in usage, a finding invisible to frequency-based approaches alone.